\section{Discussion}
The result in \autoref{fig:cosine}--\ref{fig:cosinepulse-RA} is overall in agreement with the theory presented above. It was shown that both the Leapfrog and the Upwind scheme could be used to model the advection equation, if the Courant number was chosen $\mu \leq 1$. When $\mu=1.1$, an unstable result was obtained; however, one must note that the Leapfrog scheme blew up before the Upwind scheme. Furthermore, when the Courant number was chosen as $\mu=1$, the Upwind scheme showed the smallest error. However, for $\mu=0.9$, a smaller error could be seen in the Leapfrog scheme. Furthermore, a large amplitude diffusion could be seen in the Upwind scheme, which was expected from the theory of the scheme. However the Leapfrog scheme showcased no such error, but both schemes showed an equal phase error.

The reason for the Upwind scheme showing larger errors for $\mu \neq 1$ might be due to the truncation error being larger in the Upwind scheme. Thus, one would expect larger deviations from the analytical solution in the Upwind scheme. One can then question why the Upwind scheme showed a much smaller error when the Courant number was chosen as $\mu=1$. However, one has to remember that the Leapfrog scheme is initialized with the forward Euler scheme. Thus, the initial conditions for the Leapfrog scheme may have an error from the forward Euler scheme, which might explain the larger error for the $\mu=1$ case compared to the Upwind scheme.

Although it was clear that a Courant number $\mu=1$ gave the cleanest result, one should note that this parameter will be hard to use in more difficult equations. More advanced equations will give larger computational modes, which will lead to one being reliant on filtering the data. Since one needs a non-zero $\gamma$-parameter, one cannot choose $\mu=1$. It is therefore more interesting to study what $\mu=0.9$ does to the system.

A Courant number $\mu=0.9$ proved to give stable, reliable results for the cosine wave, both for the Leapfrog scheme and the Upwind scheme. For the cosine pulse, a Courant number $\mu=0.9$ gave oscillatory results for the Leapfrog scheme and diffused results for the Upwind scheme. However, introducing the Robert-Asselin filter (with $\gamma > 0.01$) proved to remove the computational mode for the Leapfrog scheme. To reduce diffusion from the Robert-Asselin filter, a filter parameter of $\gamma=0.05$ gave the cleanest result.